
\subsection{Empirical concentration by Stone--Weierstrass approximation}
\label{leanproof:ANT-2-1}

The formal proof follows the manuscript's Stone--Weierstrass
approximation and Cauchy--Schwarz height transfer. The approximation is
proved for the actual compact cube, then extended by its coordinate clamp.

\begin{proof}[Complete proof following the Lean argument]
Fix $d\ge2$, put $p=(d-1)/d$, and let $\nu^\ast$ be the law of
$\vF^\ast$ on $E$. Write $T_n=\sum_{k=1}^n(H^{(k)})^2$.
The Gaussian power integral identifies
$\rho^\ast=\psi_E^p/\int_E\psi_E^p$.
For a bounded continuous function $b:E\to\R$, define
\[
 m_b=\int_E b\,\dd\nu^\ast,
 \qquad
 \Lambda_b(t)=\frac1p\log\int_E e^{ptb}\,\dd\nu^\ast.
\]
Subtracting the instance $(\sigma,c)=(0,0)$ of
Proposition~\ref{prop:volume} from its instance $(tb,c)$ gives the actual
conditional identity and limit
\[
 \E_n e^{t\sum_{k=1}^n b(\vF^{(k)})+cT_n}
   =\frac{\cZ_n(tb,c)}{\cZ_n(0,0)},
 \qquad
 \frac1n\log\E_n
 e^{t\sum_{k=1}^n b(\vF^{(k)})+cT_n}
 \longrightarrow\Lambda_b(t),                 \tag{P1}
\]
for every real $t$ and every $c<1/2$. These conditional expectations are
well defined: the antichain event has positive probability, and boundedness
of $b$ together with $c<1/2$ makes the weighted Gaussian integrals finite.

We first deduce two elementary exponential tails from (P1). Setting $b=0$
and taking, for example, $c=1/4$, gives
$n^{-1}\log\E_n e^{cT_n}\to0$. Consequently, for every $a>0$ and all
sufficiently large $n$,
\[
 \P_n(T_n\ge an)
 \le e^{-can}\E_n e^{cT_n}
 \le e^{-can/2}.                              \tag{P2}
\]
For the floor average, put $Z=b(\vF^\ast)-m_b$, so $\E Z=0$. Since $Z$ is
bounded, all its exponential moments are finite. Differentiation under
the integral is justified on each bounded interval of $t$ by the
uniform bounds on $Z^j e^{tZ}$ for $j=0,1,2$. Thus $M(t)=\E e^{tZ}$
and $K(t)=\log M(t)$ are twice continuously differentiable near zero.
Their values at zero are $M(0)=1$ and $K(0)=0$, and
$M'(0)=K'(0)=\E Z=0$. Taylor's theorem gives
$M(t)=1+O(t^2)$ and $K(t)=O(t^2)$, hence
$\Lambda_{b-m_b}(t)=p^{-1}K(pt)=O(t^2)$.
Given $a>0$, choose $t>0$ sufficiently small that
$\Lambda_{b-m_b}(t)<ta/2$. Markov's inequality and (P1), at $c=0$, imply
\[
 \P_n\left(\frac1n\sum_k b(\vF^{(k)})-m_b\ge a\right)
 \le \exp\{-tan+n\Lambda_{b-m_b}(t)+o(n)\}
 \le e^{-tan/4}
\]
for all sufficiently large $n$. Apply the same argument to $-b$ and
absorb the factor $2$ into a smaller positive exponential rate. We have
proved that, for every bounded continuous $b$ and every $a>0$, there is
$\alpha_b>0$ such that
\[
 \P_n\left(\left|\frac1n\sum_k b(\vF^{(k)})-m_b\right|\ge a\right)
 \le e^{-\alpha_b n}                          \tag{P3}
\]
for all sufficiently large $n$.

Now let $g$ be any continuous test on the closed cube, and fix an
approximation tolerance $\eta>0$.
Here is the precise Stone--Weierstrass step used in Lean. On a compact
metric space $K$, let $\mathcal L\subset C(K,\mathbb R)$ be the set of
Lipschitz functions, allowing an arbitrary finite Lipschitz constant.
It contains constants and is closed under pointwise minimum and maximum:
if $f,g$ have constants $L_f,L_g$, then $\min(f,g)$ and $\max(f,g)$ have
constant $\max(L_f,L_g)$.
It also separates points strongly. Given any prescribed values $v(x),v(y)$
at distinct points $x,y$, the function
\[
 z\longmapsto v(x)+\frac{v(y)-v(x)}{\operatorname{dist}(x,y)}
                                  \operatorname{dist}(x,z)
                                                        \tag{P4}
\]
is Lipschitz and takes those values; if $x=y$, use the constant function.
The lattice form of the Stone--Weierstrass theorem therefore gives
$\overline{\mathcal L}=C(K,\mathbb R)$ in the uniform norm.
Apply it to $K=[0,1]^d$: for every $\eta>0$ there is a Lipschitz function
$f$ on the cube such that
\[
                 \sup_{x\in[0,1]^d}|g(x)-f(x)|<\eta.       \tag{P5}
\]
The coordinate clamp $c(x)_i=\max(0,\min(1,x_i))$ is $1$-Lipschitz
in the coordinate maximum metric and fixes the cube. Replacing $f$ by
$f\circ c$ makes it a globally bounded Lipschitz function without
changing any cube value. It is bounded because the original $f$ is
continuous on the compact cube. A Lipschitz constant in the coordinate
maximum metric is also valid in Euclidean distance, which dominates that
metric. Thus (P5) supplies an actual bounded ambient function and its
finite Lipschitz constant; this approximation is not an assumed input.

Let $f$ satisfy (P5), with Lipschitz constant $K$.
The coordinatewise normal CDF is $L$-Lipschitz in Euclidean distance,
where $L=(2\pi)^{-1/2}$, because $0<\Phi'\le L$.
Put $b=f\circ\Phi$ on $E$. This is a bounded continuous floor test,
and $m_b=\int f\,\dd\mu^\ast$ by the actual pushforward law.
For $\vG=\vF+H\ve$ with $\|\ve\|=1$,
\[
 |f(\Phi(\vG))-f(\Phi(\vF))|\le KL|H|.
\]
For $n>0$, the triangle inequality and Cauchy--Schwarz give exactly
\[
 \left|\frac1n\sum_k f(\vX^{(k)})
       -\frac1n\sum_k f(\Phi(\vF^{(k)}))\right|
 \le\frac{KL}{n}\sum_k|H^{(k)}|
 \le KL\sqrt{\frac{T_n}{n}}.                 \tag{P6}
\]
The approximation changes each of the empirical and limiting averages
by at most $\eta$. Combining (P5)--(P6),
\[
 \left|\frac1n\sum_k g(\vX^{(k)})-\int g\,\dd\mu^\ast\right|
 \le2\eta+KL\sqrt{\frac{T_n}{n}}
       +\left|\frac1n\sum_k b(\vF^{(k)})-m_b\right|. \tag{P7}
\]
For a prescribed deviation $\varepsilon>0$, choose
$\eta=\varepsilon/8$, put $C=KL\ge0$, and set
$a_0=(\varepsilon/[4(C+1)])^2>0$.
If $T_n/n<a_0$ and the floor-test deviation is less than
$\varepsilon/4$, then the right side of (P7) is less than
$\varepsilon$. The height tail (P2), the floor tail (P3), and a union
bound therefore give an exponential tail for every fixed positive
deviation of every continuous cube test.

Finally, the weak topology on probability measures on the compact cube is
generated by the integrals of continuous functions. Any prescribed open
weak neighborhood $\mathcal U$ of $\mu^\ast$ therefore contains a set
\[
 \left\{\nu:
   \left|\int g_j\,\dd\nu-\int g_j\,\dd\mu^\ast\right|
      <\varepsilon_j\quad(1\le j\le m)\right\},
\]
for finitely many continuous tests and positive tolerances. The preceding
test bounds and a finite union bound prove the first assertion of
Theorem~\ref{thm:empirical-limit}. Explicitly, if finitely many bad events
have probabilities at most $e^{-a_jn}$ eventually, put
$a=\min_j a_j>0$. Their union has probability at most
$m e^{-an}\le e^{-an/2}$ once $n\ge2\log m/a$.
An empty collection of bad events has probability zero.

For a Borel set $B$ with $\mu^\ast(\partial B)=0$, the
Portmanteau continuity-set theorem says that
$\nu\mapsto\nu(B)$ is continuous at $\mu^\ast$.
For every positive tolerance, some open weak neighborhood of $\mu^\ast$
therefore forces that tolerance for $\nu(B)$.
Applying the first assertion gives convergence of the literal empirical
frequency of $B$ in probability, and indeed gives an exponential tail
for each fixed positive tolerance. All measures here are the actual
closed-cube measures, so no extension of the state space changes either
conclusion.
\end{proof}

\subsection{The support lemma by the written cube-shift argument}
\label{leanproof:ANT-2-2}

The support and closure formula and the common cube-coordinate shift in
Lemma~\ref{lem:support} are now proved by the argument written in the
manuscript. The earlier inverse-coordinate-centering construction is
retained only as a supplementary proof.

\begin{proof}[Complete proof following the Lean argument]
Put $S=\Phi(E)$ and
\[
 C=S\cup\{x\in[0,1]^d:\text{$x$ has a zero coordinate and a one coordinate}\}.
\]
The actual limiting measure is the pushforward of a full-support Gaussian
measure by the continuous map
$Y\mapsto\Phi(\sqrt{d/(d-1)}\Proj_EY)$.
Its range is exactly $S$: the inverse-normal coordinates always sum to zero,
and each $f\in E$ is obtained by taking $Y=f/\sqrt{d/(d-1)}$.
The support of a continuous pushforward of a full-support measure is the
closure of its range. Indeed, an open neighborhood of a range point has
a nonempty open preimage, of positive measure, and the pushforward is
concentrated on the closed closure of its range. Consequently
$\operatorname{supp}\mu^\ast=\overline S$.

We characterize that closure directly. Suppose points of $S$ approach
$x$ in the closed cube. If every coordinate of $x$ is interior, continuity
of the inverse-normal coordinates implies
$\sum_i\Phi^{-1}(x_i)=0$, so $x\in S$.
If some $x_j=0$ and no coordinate is $1$, then all other inverse-normal
coordinates are eventually bounded above, while the $j$th tends to
$-\infty$. Thus their sum tends to $-\infty$, contradicting its constant
value zero along $S$.
The case of a coordinate equal to $1$ and no zero coordinate is identical
with signs reversed: all other terms are eventually bounded below, while
one tends to $+\infty$.
This proves $\overline S\subseteq C$, including every boundary case.
The same reasoning can be expressed by limits along the neighborhood
filter restricted to $S$, so no assumption that a chosen sequence exists
is needed.

Conversely, a mixed zero--one boundary point belongs to the closure by the
explicit construction in the manuscript. Let
$I_0=\{i:x_i=0\}$ and $I_1=\{i:x_i=1\}$, both nonempty, and put
$A=\sum_{i\notin I_0\cup I_1}\Phi^{-1}(x_i)$.
For $t\to+\infty$, set
\[
 f_i(t)=
 \begin{cases}
  -t/|I_0|,&i\in I_0,\\
  (t-A)/|I_1|,&i\in I_1,\\
  \Phi^{-1}(x_i),&i\notin I_0\cup I_1.
 \end{cases}
\]
The sum is $-t+(t-A)+A=0$, and hence $\Phi(f(t))\in S$.
The endpoint limits of $\Phi$ give $\Phi(f(t))\to x$.
Together with $S\subseteq\overline S$, this proves $\overline S=C$.

Now let $x\in[0,1]^d\setminus\overline S$.
Write $m=\min_i x_i$ and $M=\max_i x_i$.
The closure formula implies $M-m<1$, because equality would give both a
zero and a one coordinate. The interval $J=(-m,1-M)$ is therefore nonempty.
For $t\in J$, all $x_i+t$ belong to $(0,1)$ and
\[
 A_x(t)=\sum_{i=1}^d\Phi^{-1}(x_i+t)
\]
is continuous and strictly increasing.
As $t\downarrow-m$, a coordinate attaining $m$ has inverse-normal value
tending to $-\infty$; all other terms are bounded above, since their upper
limits are at most $M-m<1$. Thus $A_x(t)\to-\infty$.
As $t\uparrow1-M$, a coordinate attaining $M$ has inverse-normal value
tending to $+\infty$; every other term is bounded below, since its lower
limit is at least $1-(M-m)>0$. Thus $A_x(t)\to+\infty$.
The intermediate value theorem supplies $t_0\in J$ with $A_x(t_0)=0$.
If $t_0=0$, then $x$ itself would be an interior zero-sum inverse-normal
point and therefore belong to $S$, contrary to the hypothesis.
Consequently $t_0\ne0$, and
$y_i=x_i+t_0$ lies in $S$. A positive $t_0$ gives $x\prec y$; a negative
$t_0$ gives $y\prec x$.
This proves every assertion of Lemma~\ref{lem:support} by its written route.
\end{proof}

\subsection{The rank limit by the written Hoeffding argument}
\label{leanproof:ANT-6-1}

The proof below follows the revised manuscript: condition the sorted-value
law, apply Hoeffding separately to each order statistic, take the finite
union bound, and transfer continuous tests and Hausdorff distance. The
previously proved sharp DKW theorem remains a separate library result.

\begin{proof}[Complete proof following the Lean argument]
We first justify the exact conditional laws. For a scalar iid continuous
sample, ties have probability zero. Let $\rho$ be the assignment of the
ordered values to labels and $V=(U_{(1)},\ldots,U_{(n)})$ its sorted list.
For any Borel set $A$ of sorted lists, relabeling the iid sample preserves
its law and $V$, and carries each rank assignment to every other one.
The $n!$ disjoint rank assignments therefore each have joint probability
$\mathbb P(V\in A)/n!$ with $V\in A$.
Thus the assignment is uniform and independent of the sorted list.
Apply this in each independent coordinate. The entire rank tuple is
independent of the full collection of sorted coordinate values.
Sorting the sample by its first coordinate reparameterizes the uniform
rank tuple into its common label order and its $d-1$ relative permutations.
The antichain event depends only on the relative permutations and is
exactly membership in $\mathcal C_n$.
Consequently, after conditioning, that tuple is uniform on $\mathcal C_n$
and the joint law of all sorted coordinate values is unchanged.
This also gives $|\mathcal C_n|=(n!)^{d-1}a_n$.
The conditional law is legitimate because $a_n>0$ for $d\ge2$.

Fix a coordinate, $n\ge1$, and $t\in[0,1]$. Sorting preserves the number
of sample values at most $t$. Before conditioning that count is a sum
of $n$ independent Bernoulli variables with success probability $t$.
More explicitly, each prescribed set of $k$ successful labels has
probability $t^k(1-t)^{n-k}$; the $\binom nk$ such disjoint choices give
\[
 \mathbb P_n(N_i(t)=k)=\binom nk t^k(1-t)^{n-k},
 \qquad 0\le k\le n.                           \tag{H1}
\]
The unchanged sorted law proves (H1) under $\mathbb P_n$ as well,
including $t=0,1$ with the usual $0^0=1$ convention.

For completeness, the Hoeffding bound used in Lean follows from the
bounded-variable subgaussian estimate. If $B$ is Bernoulli with mean $t$,
then
\[
        \mathbb E e^{\lambda(B-t)}\le e^{\lambda^2/8}
        \qquad(\lambda\in\mathbb R).
\]
One can see the constant directly: the log moment-generating function
has value and derivative zero at zero; its second derivative is the
variance of a tilted Bernoulli variable, which is at most $1/4$.
Integrating that derivative bound twice gives the displayed inequality.
At $t=0,1$ the centered variable is identically zero.
Independence and exponential Markov, with $\lambda=4a/n$, therefore give
both tails
\[
 \mathbb P_n(N_i(t)-nt\ge a)\le e^{-2a^2/n},\qquad
 \mathbb P_n(nt-N_i(t)\ge a)\le e^{-2a^2/n}
 \quad(a>0).                                  \tag{H2}
\]
The formal proof uses the same unit-width Bernoulli range and its
subgaussian parameter $1/4$, rather than a DKW premise.

Write $u_{r,n}=(r-1/2)/n$ for $1\le r\le n$, and fix $s>0$.
If $S_{i,r}>u_{r,n}+s$ and $t=u_{r,n}+s\le1$, then
$N_i(t)\le r-1$, whereas $nt=r-1/2+ns$.
This is a downward deviation of at least $ns$.
The threshold is positive; if it exceeds $1$, the event is empty.
Similarly, if $S_{i,r}<u_{r,n}-s$ and $t=u_{r,n}-s\ge0$, then
$N_i(t)\ge r$ and $nt=r-1/2-ns$, an upward deviation of at least $ns$.
If this threshold is negative, the event is empty.
Apply (H2) with $a=ns$ to obtain the exact finite-sample bound
\[
 \mathbb P_n\bigl(|S_{i,r}-u_{r,n}|>s\bigr)
                       \le2e^{-2ns^2}.        \tag{H3}
\]
The two endpoint threshold cases $t=0,1$ are included.

Relabel the points in increasing first-coordinate order as
$\widetilde{\mathbf X}^{(k)}$. Its $i$th coordinate is
$S_{i,\pi_i(k)}$, while the corresponding midpoint rank coordinate is
$u_{\pi_i(k),n}$. Thus
\[
 D_n:=\max_k\|\widetilde{\mathbf X}^{(k)}-\mathbf R^{(k)}\|
       \le\sqrt d\max_{i,r}|S_{i,r}-u_{r,n}|.
\]
A common relabeling changes neither maximum nor empirical measure.
Apply (H3) at tolerance $s/\sqrt d$ and take the union over the $dn$
coordinate/rank choices. This proves exactly
\[
       \mathbb P_n(D_n>s)\le2dn e^{-2ns^2/d}
       \qquad(n\ge1,\ s>0).                  \tag{H4}
\]
For a fixed $s>0$, the prefactor is absorbed into a smaller exponential
rate: $\log(2dn)/n\to0$, so, eventually,
$2dn e^{-2ns^2/d}\le e^{-ns^2/d}$.

Let $\mathcal U$ be any prescribed open weak neighborhood of $\mu^\ast$
in probability measures on the closed cube. Choose finitely many
continuous tests $g_j$ and $\varepsilon>0$ so that all integral
deviations less than $\varepsilon$ imply membership in $\mathcal U$.
Uniform continuity supplies $s>0$ such that moving any cube point by at
most $s$ changes every $g_j$ by less than $\varepsilon/3$.
On $\{D_n\le s\}$, replacing original points by their rank midpoints
therefore changes each empirical test average by less than
$\varepsilon/3$.
Spreading a midpoint over its rank box moves it by at most
$\sqrt d/(2n)$, and hence changes every test average by less than another
$\varepsilon/3$ for large $n$.
Theorem~\ref{thm:empirical-limit} at the finitely many tolerances
$\varepsilon/3$, together with (H4), proves simultaneous membership of
both rank measures in $\mathcal U$ outside a finite union of exponentially
rare events. If those probabilities are at most $e^{-a_j n}$ eventually,
choose $a=\min_j a_j>0$ and use $m e^{-an}\le e^{-an/2}$ for large $n$.
An empty test family has no bad event.
This proves Theorem~\ref{thm:permuton}.

The rank-box measure has exactly uniform coordinate marginals: every
coordinate rank occurs once, and the equal mixture of the $n$ uniform
interval laws on $[(r-1)/n,r/n]$ is uniform measure on $[0,1]$.
The shared endpoints have measure zero. The midpoint marginal is exactly
uniform on its midpoint grid.

Finally, two finite clouds matched by their labels satisfy
\[
 d_{\rm H}(\{\widetilde{\mathbf X}^{(k)}\}_k,
                    \{\mathbf R^{(k)}\}_k)\le D_n.
\]
The Hausdorff triangle inequality gives
\[
 d_{\rm H}(A_n^{\mathrm{rk}},\overline{\mathcal S_d})
             \le D_n+d_{\rm H}(A_n,\overline{\mathcal S_d}).
\]
For any fixed $\delta>0$, (H4) at $s=\delta/2$ and
Theorem~\ref{thm:shape} at $\delta/2$ show that a distance greater than
$\delta$ is exponentially rare. Absorbing their two-event union into a
smaller positive rate proves Remark~\ref{rem:rank-hausdorff}.
\end{proof}
